\section{Economics and Reproducibility}
\label{sec:recovery}
\sectionaim{Determine when recovery repays its information cost and whether the decision chain can be replayed.}

\subsection{RQ8: Is Recovery Economically Worthwhile?}
\textbf{Question.} When do target profiles repay their acquisition cost? \textbf{Protocol.} In the registered distance-count model,
\begin{align}
 C_\pi(N)&=C_{\mathrm{acquire},\pi}+N\bar c_\pi,\label{eq:cost}\\
 N_{\mathrm{BE}}&=\frac{C_{\mathrm{acquire},\pi}-C_{\mathrm{acquire},\mathrm{ref}}}
 {\bar c_{\mathrm{ref}}-\bar c_\pi},\label{eq:breakeven}
\end{align}
with separate handling when serving savings are nonpositive. Target truth, selection, certification, and search are charged; source-profile truth is treated as shared or previously acquired. \textbf{Observation.} Dataset-level break-even is \WZeroSiftBreakEven\ query executions on SIFT and \WZeroArxivBreakEven\ on Arxiv. All SIFT targets amortize; one Arxiv target has no positive serving saving. At one million executions, modeled average savings are \WZeroSiftNetSavingMillion\ and \WZeroArxivNetSavingMillion\ distance evaluations per target. \textbf{Conclusion.} TCP has positive modeled lifecycle value for sufficiently long-lived recurring-query workloads, while short lifecycles and nonamortizing targets should retain the endpoint. The economic result is conditional on symmetric treatment of previously acquired source profiles.

\begin{table}[t]
\caption{Lifecycle summary. Break-even counts query executions; net savings are modeled distance evaluations per target at $N=10^6$.}
\label{tab:economics}
\small
\begin{tabular*}{\linewidth}{@{\extracolsep{\fill}}lrrr@{}}
\toprule
Dataset & Break-even & Nonamortizing & Net saving at $10^6$ \\
\midrule
SIFT & \WZeroSiftBreakEven & \WZeroSiftNonamortizing & \WZeroSiftNetSavingMillion \\
Arxiv & \WZeroArxivBreakEven & \WZeroArxivNonamortizing & \WZeroArxivNetSavingMillion \\
\bottomrule
\end{tabular*}
\end{table}

\subsection{RQ9: Can the Evidence and Decision Be Reproduced?}
\textbf{Question.} Can an artifact evaluator regenerate the principal numbers and inspect every decision boundary? \textbf{Protocol.} The artifact freezes input hashes, role counts, implementation scripts, per-build outputs, decision manifests, and a single evidence ledger. A deterministic checker reselects unique rows, recomputes risks and gains, verifies per-build aggregation, and regenerates all headline LaTeX macros. The replay protocol preserves source/target IDs and selection/certification/evaluation separation; portable smoke checks cover table regeneration before full replay. \textbf{Observation.} The W4 manuscript uses 57 ledger-bound frozen values, and the checker rejects input drift or inconsistent aggregation. All principal tables are generated from committed CSV/JSON rather than transcribed from logs. \textbf{Conclusion.} The numerical claims and audit chain are mechanically reviewable from the frozen artifact. This reproducibility guarantee covers committed summaries and code paths; independent raw-query replication remains the artifact-evaluation task.

\paragraph{Evidence hierarchy.}
RQ1, RQ4, and RQ5 provide the decisive result: direct reuse violates the target while isolated recalibration recovers empirical safety and efficiency. RQ6--RQ8 establish tail, robustness, and economic conditions. RQ2--RQ3 extend the phenomenon across methods, graph families, and scale. RQ9 binds these claims to replayable evidence without collapsing protocol differences into a single leaderboard.
